% TOP_PROBLEM: {"problemId": "problem.minkowskis-conjecture-on-products-of-non-homogeneous-linear-forms", "canonicalTitle": "Minkowski's conjecture on products of non-homogeneous linear forms", "releaseRank": 429, "primaryDomain": "number_theory_arithmetic_geometry", "primaryDomainLabel": "Number theory & arithmetic geometry", "displayStatus": "open_with_solved_subcases", "releaseStatus": "open_with_solved_subcases"}
% !TeX program = lualatex
% ATTEMPT_STATUS: unresolved
\documentclass[11pt]{article}
\usepackage[margin=25mm]{geometry}
\usepackage{fontspec}
\setmainfont{Latin Modern Roman}
\usepackage{amsmath,amssymb,mathtools}
\usepackage{unicode-math}
\setmathfont{Latin Modern Math}
\usepackage{xurl,hyperref}
\hypersetup{colorlinks=true,urlcolor=blue}
\setcounter{secnumdepth}{0}
\setlength{\parindent}{0pt}
\setlength{\parskip}{5pt}
\setlength{\emergencystretch}{3em}
\begin{document}
\section*{\texorpdfstring{Minkowski's conjecture on products of non-homogeneous linear forms}{Minkowski's conjecture on products of non-homogeneous linear forms}}\label{429}
\subsection{Definitions and mathematical statement}
Let $A\in\operatorname{GL}_n({\mathbb{R}})$ satisfy $|\det A|=1$, so that $\Lambda=A{\mathbb{Z}}^n$ is a lattice of covolume one. For each translate $x+\Lambda$, the assertion is that an actual lattice translate point lies in the closed multiplicative region
\[
 \exists m\in{\mathbb{Z}}^n:\quad \prod_{i=1}^n|x_i+(Am)_i|\le2^{-n}.
\]
The quantifiers are every integer $n\ge1$, every such matrix, and every $x\in{\mathbb{R}}^n$. For a lattice of arbitrary covolume $D>0$ the corresponding bound is $D/2^n$. A statement about the infimum of these products being at most the threshold does not, without an attainment argument, supply the requested point at the threshold.
\par\smallskip\textit{Formulation and concept references:} \href{https://arxiv.org/pdf/2009.09992v1}{[S1]}.

\subsection{Short English statement}
Does every translate of a lattice contain a point whose coordinate product is at most its covolume divided by two to the dimension?
\subsection{Research attempt}
\paragraph{Process and approach.}
This attempt was generated by GPT-6 Astra Ultra on 27 September 2026. We prove the desired constant for lattices with a triangular basis in the fixed coordinate system, extend this to rational matrices by integral column operations, and investigate whether density of rational matrices finishes the problem. It does not: the argument gives no uniform bound on the integer witnesses during approximation. We also construct an example showing why the coordinate-product objective cannot be treated as a proper distance function.

The inherited source [S1] proves a finite-dimensional case through lattice reduction and a covering problem. Its introduction explicitly distinguishes Minkowski's conjecture from the stronger Woods covering conjecture; failure of the latter in other dimensions would not disprove the former. Our argument is independent of those finite-dimensional estimates and makes no claim to settle the arbitrary-dimensional problem.

\paragraph{Normalization and transformations that preserve the objective.}
It is convenient to allow $D=|\det A|>0$ and prove the bound $D/2^n$. Replacing $A,x$ by $D^{-1/n}A,D^{-1/n}x$ shows equivalence with the determinant-one normalization.

Right multiplication of $A$ by a unimodular integer matrix $U$ does not change its lattice: $AU\mathbb Z^n=A\mathbb Z^n$. Permuting output coordinates also preserves the absolute product. Multiplying coordinate $i$ by a nonzero factor $b_i$ multiplies both the coordinate product and the covolume by $\prod_i|b_i|$. These operations therefore preserve the assertion. Arbitrary rotations do not preserve a coordinate product, a distinction that matters below.

\paragraph{A triangular rounding lemma in every dimension.}
Suppose $A=(a_{ij})$ is upper triangular with nonzero diagonal. Fix $x\in\mathbb R^n$. Choose $m_n\in\mathbb Z$ nearest to $-x_n/a_{nn}$. Then
\[
 |x_n+a_{nn}m_n|\le |a_{nn}|/2.
\]
After $m_{i+1},\ldots,m_n$ have been selected, choose $m_i$ nearest to
\[
 -\frac{x_i+\sum_{j>i}a_{ij}m_j}{a_{ii}}.
\]
This choice makes $|x_i+(Am)_i|\le|a_{ii}|/2$ and does not change any previously treated coordinate, since the entries below the diagonal vanish. Induction therefore gives an actual integer vector $m$, with
\[
 \prod_{i=1}^n|x_i+(Am)_i|
 \le2^{-n}\prod_i|a_{ii}|=2^{-n}|\det A|.
 \tag{429.1}
\]
Ties in nearest-integer rounding can be settled either way, and negative diagonal entries cause no change in the absolute inequalities. Lower triangular matrices are treated by choosing coordinates in the opposite order. This is a constructive proof, not merely an estimate of an infimum.

The constant is sharp even in the diagonal subclass. For $A=\operatorname{diag}(a_1,\ldots,a_n)$ and $x=(a_1/2,\ldots,a_n/2)$, every integer vector satisfies
\[
 |x_i+(Am)_i|=|a_i|\,|m_i+1/2|\ge|a_i|/2.
\]
Their product is at least $D/2^n$, and equality is attained by choosing each $m_i$ to be $0$ or $-1$. Hence a universal smaller constant is impossible.

\paragraph{A block extension.}
Suppose $A$ is upper block triangular, with invertible diagonal blocks $A_1,\ldots,A_r$ of sizes $n_1,\ldots,n_r$. Assume the desired inequality is known for each of these particular blocks, for every translate. First select the last block of integer coordinates using the assertion for $A_r$. Its contribution to the coordinate product is at most $|\det A_r|/2^{n_r}$. The chosen coordinates shift the translation in each earlier block but do not alter its matrix. Proceed upward through the blocks.

Multiplying the inequalities yields $|\det A|/2^n$, because $n=\sum_jn_j$ and the determinant of a block triangular matrix is the product of the diagonal block determinants. This closure argument shows exactly how established lower-dimensional cases can be combined when the required block structure is already present in the fixed coordinates. It supplies no such block structure for a general real matrix.

\paragraph{All rational matrices admit the needed column reduction.}
Let $A$ have rational entries and nonzero determinant. Choose a positive integer $q$ such that $B=qA$ is integral. There is a unimodular integer matrix $U$ for which $BU$ is upper triangular. Here is an explicit inductive reason, avoiding a change of the output coordinate system.

The last row of $B$ is a nonzero integer row. The Euclidean algorithm, implemented by swaps, sign changes, and integer additions of columns, transforms that row to $(0,\ldots,0,g)$, where $g>0$ is the greatest common divisor of its entries. Every such column operation is unimodular. The resulting matrix has the form
\[
 \begin{pmatrix}C&b\\0&g\end{pmatrix},
\]
with $C$ an integral invertible $(n-1)$ by $(n-1)$ matrix because the full determinant is nonzero. Apply the same argument to $C$, using only the first $n-1$ columns. These operations preserve the zeros in the last row. Induction proves the triangularization assertion.

Consequently $AU=q^{-1}BU$ is upper triangular and generates the same lattice as $A$. Applying (429.1) to $AU$ and returning to the original integer coordinates proves the conjectured inequality for every rational $A$ and every real translate $x$. The same proof covers a nonzero real scalar multiple of a rational matrix.

\paragraph{Why density leaves an unproved compactness step.}
Choose nonsingular rational matrices $A_j\to A$. For every $j$, the preceding construction gives some $m_j\in\mathbb Z^n$ satisfying
\[
 \prod_i|x_i+(A_jm_j)_i|\le2^{-n}|\det A_j|. \tag{429.2}
\]
If the sequence $m_j$ had a bounded subsequence, discreteness would give a further subsequence with constant value $m$. Passing to the limit in (429.2) would then prove the desired inequality for $A$. This establishes a precise conditional reduction.

But the construction does not bound $m_j$. The integral column operations may involve increasingly large denominators and basis coefficients. More fundamentally, a bound on a product of coordinates does not bound the vector itself: one coordinate can approach zero while another tends to infinity. Thus (429.2) provides no compact region from which to extract the required constant subsequence. Density of the solved family alone is insufficient. This identifies a missing argument, not a proof that all possible choices of witnesses must escape.

Using real orthogonal triangularization is a different invalid shortcut. If $A=QR$ with $R$ triangular, rotating $x$ by $Q^{-1}$ changes the coordinate product. For example $(L,0)$ has product zero in dimension two, whereas its rotation by $45$ degrees has absolute coordinate product $L^2/2$. Equal Euclidean lengths and equal determinants do not fix this discrepancy.

\paragraph{An explicit nonattained infimum illustrating nonproperness.}
Let
\[
 \alpha=2\sum_{j=1}^{\infty}3^{-j!},\qquad
 A=\begin{pmatrix}1&0\\\alpha&1\end{pmatrix},\qquad
 x=(1/2,\alpha/2).
\]
Write the $N$th partial sum of $\alpha$ as $p_N/q_N$, where $q_N=3^{N!}$ and $p_N=2\sum_{j\le N}3^{N!-j!}$ is even. The denominators are odd, and the positive tail satisfies
\[
 0<|\alpha-p_N/q_N|\le C3^{-(N+1)!}.
\]
This also proves $\alpha$ irrational: a distinct rational with fixed denominator would differ from $p_N/q_N$ by at least a constant times $q_N^{-1}$, contradicting the displayed tail estimate for large $N$.

Take $m_1=(q_N-1)/2$ and $m_2=-p_N/2$. Then
\[
 x+Am=(q_N/2,(q_N\alpha-p_N)/2),\qquad
 \prod_{i=1}^2|x_i+(Am)_i|
 \le C3^{2N!-(N+1)!}\longrightarrow0.
 \tag{429.3}
\]
For every integer pair, however, the first coordinate is a nonzero half-integer and the second is $\alpha(m_1+1/2)+m_2\ne0$ by irrationality. The product has infimum zero and never attains zero. This example itself satisfies the target by triangular rounding. In particular, an infimum strictly smaller than $1/4$ already ensures some value below $1/4$. The example only disproves the unjustified general assertion that the infimum of the product must be attained; it is not a counterexample to Minkowski's bound.

\paragraph{Verification and outcome.}
Exact rational checks verify the rounding construction on selected triangular matrices and translates, the sharp diagonal examples, and the integer parity in (429.3). The general Euclidean-algorithm proof and the limiting estimates are given above rather than inferred from samples.

\textbf{UNRESOLVED.} The desired sharp constant is established for triangular bases, the stated block extensions, and rational matrices. Extending by density stops at an unproved bound on the integer witnesses; no argument closes that gap for arbitrary real lattices in every dimension.

\subsection{Sources}
\begin{itemize}
\item[S1]Leetika Kathuria and Madhu Raka, On conjectures of Minkowski and Woods for n=10, arXiv2009.09992v1.
\url{https://arxiv.org/pdf/2009.09992v1}.
\textit{Formulation source.}
\end{itemize}
\end{document}
